% Appendix: u-adic convergence; the statement and proof of lem:tra-conv, used in sections 4, 5 and 7.

\section{$u$-adic convergence}
\label{sec:tra-convproof}

This appendix states and proves \cref{lem:tra-conv}, which justifies the infinite sums of \cref{sec:tra,sec:red,sec:flh}; the notation is that of \cref{sec:tra-tower}.

\begin{lem}[$u$-adic convergence]\genusstatus{AI-proposed, unchecked}
\label{lem:tra-conv}
Work in the $u$-adic setting. A flat model correlator is an ancestor correlator $\langle z^{k_1}H^{a_1}\cdots z^{k_n}H^{a_n}\rangle_h(v)$ of \cref{def:set-modcohft}, evaluated at $v = v(u)$; it vanishes unless $\sum_ik_i\leq\mathsf D_{h,n} := 3h-3+n$, because the class of the CohFT has no component of negative degree and $\dim\overline{\mathcal M}_{h,n} = \mathsf D_{h,n}$.
\begin{description}
\item[(a) Frames] For every family $(c_{(k,a)})$ in $\bbQ\fps{u}$ the family $(c_{(k,a)}\fr(k,a))$ is summable, the $z^jH$-coefficient of $\fr(k,a)$ having $u$-order at least $(4(k-j)+\nu_a+1)/5$; the decompositions of \cref{lem:tra-frames}~(c) and \cref{lem:tra-twoslot} hold with these sums.
\item[(b) Frame expansions inside correlators] Let $h\geq 0$, $n\geq 1$ and $2h-2+n>0$.
(b1) The correlator extends to a symmetric $\bbQ\fps{u}$-multilinear map $\langle\ \rangle_h\colon(\cH^u_+)^n\to\bbQ\lser{u}$,
\[
\langle X_1\cdots X_n\rangle_h := \sum_{\sum k_i\leq\mathsf D_{h,n}}\\prod_i(X_i)_{(k_i,a_i)}\,\langle z^{k_1}H^{a_1}\cdots z^{k_n}H^{a_n}\rangle_h(v(u)),
\]
a finite sum, which is of finite type in each slot.
(b2) Let $B := \{(k,a) : a\neq 1\}\cup\{H\}$, $t_{(k,a)} := \fr(k,a)$ and $t_H := H$. If $X_i = \sum_{\beta\in B}c_{i,\beta}t_\beta$ with $c_{i,\beta}\in\bbQ\fps{u}$, as is every element of $\sK$ (\cref{lem:tra-frames,cor:tra-slice}), then the family $(\prod_ic_{i,\beta_i}\langle t_{\beta_1}\cdots t_{\beta_n}\rangle_h)_{\beta\in B^n}$ is summable in $\bbQ\lser{u}$ and its sum is $\langle X_1\cdots X_n\rangle_h$. Every term is a correlator of homogeneous elements of $\sK$ over $\bbQ[v]$ evaluated at $v(u)$; no flat component $z^jH$, $j\geq 1$, of a frame occurs as an insertion of its own.
(b3) For $y\in z^2\cH^u_+$ the vertex factor $\sum_{m\geq 0}\frac1{m!}\langle X_1\cdots X_ny^m\rangle_h$ is a finite sum ($m\leq\mathsf D_{h,n}$ by (d)) and is of finite type in each $X_i$. If moreover $y\in\sK$, write $y = y^\circ+\delta H$ with $\delta := \Pi_0(y)$ and $y^\circ\in\sKo$ (\cref{cor:tra-slice}); then $y^\circ_{(k,a)} = y_{(k,a)} = 0$ for $k\leq 1$ and $a\neq 1$, and for each $m$ the expansion of $y^m = (y^\circ+\delta H)^m$ is finite, each of its terms being covered by (b2) with $n+m$ slots.
(b4) Let $\Gamma$ be a leg-free stable graph with at least one edge, so that every vertex carries $n_w\geq1$ half-edges (\cref{lem:red-genus}). Put on each edge a bivector $E_e\in\bbH\otimes\bbH\otimes\bbQ\fps{\psi,\psi'}\fps{u}$ with property (i) of \cref{lem:tra-twoslot}, and at each vertex $w$ a vertex factor as in (b3), of genus $g_w$, with one slot per half-edge and a leaf vector $y_w\in z^2\cH^u_+\cap\sK$. Then the contraction of these data along $\Gamma$ is a finite sum, and it is the sum of the summable family obtained by expanding every $E_e$ by \cref{lem:tra-twoslot} and every $y_w$ as in (b3). Every member of the family is a product of elements of $\bbQ\fps{u}$ and of correlators of homogeneous elements of $\sK$ in which every half-edge slot carries a frame.
\item[(c) Vanishing series; pairings] The coefficient $\coef{z^{-n}}\YC_b(v(u),z)$, given in closed form by \cref{lem:gau-vanshape}, has $u$-order at least $n-\lfloor n/5\rfloor\geq 4n/5$. Hence $\YC(v(u))\in\bbQ[z^{-1}]\fps{u}^5$, $J^{\mathrm C}=\eta^{-1}\YC$ lies in $\cH^u$ and $z^{-k}J^{\mathrm C}\in\cH^u_-$ for $k\geq 1$, and every pairing and product used in \cref{sec:tra} ($\Pi_i$ on $\cH^u_+$; $P\YC$, $P^{\mathrm{vec}}J^{\mathrm C}$ and $Rf$ for matrices $P$, $R$ over $\bbQ\fps{u,z}$; $\Omega_{\mathrm s}$ on $\cH^u$) is a ring operation in $\bbQ\lser{z}\fps{u}$, i.e.\ a sum that is finite modulo every power of $u$. (The contraction of \cref{lem:tra-E} is treated in the same way, in the ring $\bbQ\lser{\psi}\fps{u,\psi'}$, in the proof of that lemma.)
\item[(d) Leaves] A flat insertion $z^kH^b$ with $k\geq 2$ has $\wt(k,b)\geq 3/5$. At a vertex of genus $h$ with $n$ slots and $m$ further insertions of $\psi$-degree $\geq 2$, a non-zero correlator has slot degrees $\sum_ik_i\leq 3h-3+n-m$; in particular $m\leq 3h-3+n$.
\item[(e) Vertex resummation] Let $g\geq 2$ and $v'(0)\neq 0$. Let $\tilde q\in -z1+z^2\cH^u_+$ with $\Pi_i(\tilde q) = 0$ for $i\geq 1$ and $\delta := \Pi_0(\tilde q)-v\in u^2\bbQ\fps{u}$, and write $y := \tilde q+z1 = y^\circ+\delta H$ with $y^\circ\in\sKo$ (\cref{cor:tra-slice}). (e1) $\langle y^m\rangle_g = 0$ for $m>3g-3$. (e2) The family $(\frac{\delta^j}{l!\,j!}\langle(y^\circ)^lH^j\rangle_g)_{l,j\geq 0}$, every member with $(l,j)\neq(0,0)$ expanded in frames by (b2), is summable, and (here and below $H^j$ under a bracket denotes $j$ insertions of $H=H^1$, not the flat class $H^j$)
\[
\overline F_g(\tilde q) = \sum_m\frac1{m!}\langle y^m\rangle_g = \sum_{l,j\geq 0}\frac{\delta^j}{l!\,j!}\langle(y^\circ)^lH^j\rangle_g,
\]
where $\langle(y^\circ)^0H^0\rangle_g := \langle\ \rangle_g = \Fm_g(v)$. (e3) If $\ML{g}$ holds (\cref{def:red-ML}), then $\sum_j\frac{\delta^j}{j!}\langle H^j\rangle_g = c_g(v+\delta)^{2-2g}$ and $\sum_j\frac{\delta^j}{j!}\sum_{l\geq 1}\frac1{l!}\langle(y^\circ)^lH^j\rangle_g\in\bbQ\fps{u}$; hence $\overline F_g(\tilde q)\in c_g\,\Pi_0(\tilde q)^{2-2g}+\bbQ\fps{u}$.
\item[(f) Dilaton re-expansion] Let $2h-2+n>0$, $X_1,\dots,X_n\in\cH^u_+$, $\varkappa\in\bbQ\fps{u}^\times$, $y\in z^2\cH^u_+$ and $c := 1-1/\varkappa$. (f1) If $\varkappa(0) = 1$, the family $(\frac{c^{m_1}}{m_1!\,m_2!}\langle X_1\cdots X_n(z1)^{m_1}y^{m_2}\rangle_h)_{m_1,m_2\geq 0}$ is summable and
\[
\sum_{m\geq 0}\frac1{m!}\langle X_1\cdots X_n(cz1+y)^m\rangle_h = \varkappa^{2h-2+n}\sum_{m\geq 0}\frac1{m!}\langle X_1\cdots X_n(\varkappa y)^m\rangle_h ,
\]
the right side being a finite sum ($m\leq 3h-3+n$). (f2) If $\varkappa(0)\neq 1$ and $\langle X_1\cdots X_ny^{m}\rangle_h\neq 0$ for some $m$, the terms of the series on the left do not tend to $0$, so it has no $u$-adic limit.
\end{description}
Beyond \cref{conv:tra-forms}, \cref{lem:tra-sum,lem:tra-frames,cor:tra-slice,lem:tra-twoslot} (which precede \cref{sec:tra-dress} and do not use this lemma), \cref{lem:gau-vanshape} in (c) and the definitions of \cref{def:set-graphsum,def:set-modcohft}, the parts depend on the following statements only: (b) and (e) on the homogeneity of \cref{lem:set-modhom}; (f) on the dilaton equation of \cref{lem:set-dilaton}; (e3) on $\ML{g}$, as a hypothesis; (e) on $v'(0)\neq 0$. In the applications, $\delta\in u^2\bbQ\fps{u}$ in (e) is \cref{lem:tra-T}~(d), and $\varkappa(0) = 1$ in (f) holds for $\Pc$ by \cref{lem:gau-diag} (\cref{prop:tra-can}).
\end{lem}

\begin{proof}
(c) By \cref{lem:gau-vanshape} (for the residue solution of \cref{def:gau-van,lem:gau-vanmod}) the $z^{-n}$-coefficient is a multiple of $v^{n-\lfloor n/5\rfloor}$, and $\ord_uv(u)\geq 1$; so modulo $u^N$ only finitely many $n$ occur, $\YC(v(u),z)\in\bbQ[z^{-1}]\fps{u}^5$, and the statements on $J^{\mathrm C}$, $z^{-k}J^{\mathrm C}$ follow. The rest is \cref{conv:tra-forms}: $\bbQ\lser{z}\fps{u}$ is a ring containing $\bbQ\fps{u,z}$ and $\bbQ[z^{-1}]\fps{u}$, in which products, coefficient extraction, the splitting, $z\mapsto -z$ and $\Res_z$ are coefficientwise in $u$ and finite modulo every power of $u$.

(a) This is the proof of \cref{lem:tra-frames}~(c) and of \cref{lem:tra-twoslot}, the $u$-order being the exponent of \cref{lem:tra-frames}~(b) with $\ord_uv\geq 1$.

(b1) Only the finitely many flat slots with $\sum k_i\leq\mathsf D_{h,n}$ contribute. By \cref{lem:set-modhom} each flat correlator is $Cv^e$ with $C\in\bbQ$ and $e = 2-2h+\sum_i\wt(k_i,a_i)\geq 2-2h-n$, as $\wt(k,a)\geq\wt(0,1) = -1$; so at $v = v(u)$ all of them lie in $u^{e_0}\bbQ\fps{u}$ with $e_0 := \min(0,(2-2h-n)\ord_uv)$. Hence $\langle X_1\cdots X_n\rangle_h\in u^{e_0}\bbQ\fps{u}$ for all $X_i\in\cH^u_+$; it is symmetric and $\bbQ\fps{u}$-multilinear, and with $X_2,\dots,X_n$ fixed it is of finite type in $X_1$: modulo $u^N$ it depends only on the $(X_1)_{(k,a)}$ with $k\leq\mathsf D_{h,n}$, taken modulo $u^{N-e_0}$.

(b2) For $\beta_i\in B$ put $\kappa(\beta_i) := k$ if $\beta_i = (k,a)$ and $\kappa(H) := 0$. By (b1) the term of $\beta = (\beta_1,\dots,\beta_n)$ is the finite sum over flat slots of $\prod_ic_{i,\beta_i}(t_{\beta_i})_{(k_i,a_i)}\langle\text{slot}\rangle_h(v(u))$. A frame $t_{\beta_i}$ with $\kappa(\beta_i)>\mathsf D_{h,n}$ enters a slot with $k_i\leq\mathsf D_{h,n}$ only through its $H$-coefficient at $z^{k_i}$, of $u$-order at least $(4(\kappa(\beta_i)-\mathsf D_{h,n})+1)/5$ by (a). So the term has $u$-order at least $e_0+\sum_i\max(0,(4(\kappa(\beta_i)-\mathsf D_{h,n})+1)/5)$, and for each $N$ only finitely many $\beta$ give a term outside $u^N\bbQ\fps{u}$: the family is summable. By \cref{lem:tra-sum} in the first slot (finite type by (b1); the family $(c_{1,\beta}t_\beta)_\beta$ is summable by (a)), $\langle X_1\cdots X_n\rangle_h = \sum_{\beta_1}c_{1,\beta_1}\langle t_{\beta_1}X_2\cdots X_n\rangle_h$; then in the second slot for each $\beta_1$, and so on. This expresses $\langle X_1\cdots X_n\rangle_h$ as an iterated sum of the family, which equals its unordered sum by regrouping. Each $t_\beta$ is a frame or $H$, a homogeneous element of $\sK$ over $\bbQ[v]$ evaluated at $v(u)$ (\cref{lem:tra-frames}~(b), (d)); the expansion is indexed by $B$, so a flat $z^jH$ with $j\geq 1$ never appears as an insertion.

(b3) By (d), $\langle X_1\cdots X_ny^m\rangle_h = 0$ for $m>\mathsf D_{h,n}$, since all components of $y$ have $\psi$-degree $\geq 2$; for each $m$ the term is multilinear, hence of finite type in each $X_i$ by (b1). If $y\in\sK$, then $y^\circ := y-\delta H\in\sKo$ by \cref{cor:tra-slice}, its frame coordinates are $y^\circ_{(k,a)} = y_{(k,a)}$ for $a\neq 1$, which vanish for $k\leq 1$ because $y\in z^2\cH^u_+$, and the multinomial expansion of $(y^\circ+\delta H)^m$ is finite.

(b4) For an edge between the vertices $w$ and $w'$ only its coefficients of $\psi^k\psi'^l$ with $k\leq\mathsf D_w := 3g_w-3+n_w$ and $l\leq\mathsf D_{w'}$ matter, by (d); so the contraction is a finite sum of products. As a function of one edge bivector, the others fixed, it is linear and, by (b1), modulo $u^N$ it depends on finitely many coefficients of that bivector, taken modulo a fixed power of $u$: it is of finite type. Expanding every edge by \cref{lem:tra-twoslot} and every leaf as in (b3) gives a family indexed by the frame labels of all half-edges and leaves; by the bound of (b2) at each vertex it is summable, and by \cref{lem:tra-sum}, applied edge by edge, its sum is the contraction. The edge coefficients $E^{\alpha\beta}_e$ lie in $\bbQ\fps{u}$, and each half-edge slot carries a frame.

(d) $\wt(k,b) = (4k+\nu_b-4)/5\geq(8-1-4)/5 = 3/5$ for $k\geq 2$, since $\nu_b\geq -1$. A correlator with $n$ slots of degrees $k_i$ and $m$ insertions of degrees $l_j\geq 2$ at genus $h$ vanishes unless $\sum k_i+\sum l_j\leq 3h-3+n+m$; so $\sum k_i\leq 3h-3+n+m-2m = 3h-3+n-m$, and $m\leq 3h-3+n$ because $\sum k_i\geq 0$.

(e) (e1) is (d) with no slot. (e2) For fixed $m$, multilinearity and symmetry give
\[
\langle y^m\rangle_g = \sum_{l+j=m}\frac{m!}{l!\,j!}\,\delta^j\langle(y^\circ)^lH^j\rangle_g .
\]
Expand $y^\circ$ in frames; only frames with $k\geq 2$ occur (b3), and they have $\wt(k,a)\geq(8+0-4)/5 = 4/5$, as $a\neq 1$ gives $\nu_a\geq 0$. By \cref{lem:set-modhom} a correlator of $l$ such frames and $j$ insertions $H$ ($\wt(0,1) = -1$) is $Cv^e$ with $e\geq 2-2g+4l/5-j$. Since $\ord_uv = 1$ (here $v'(0)\neq 0$ is used: for negative $e$, $\ord_uv^e = e\ord_uv$) and $\ord_u\delta^j\geq 2j$, every term of the expanded $(l,j)$-family has $u$-order at least $2-2g+j+4l/5$. So for each $N$ only finitely many $(l,j)$ contribute outside $u^N\bbQ\fps{u}$, and for each $(l,j)$ the frame family is summable by (b2); the family over $(l,j)$ and the frame labels is summable. Grouping it by $m = l+j$ gives $\sum_m\frac1{m!}\langle y^m\rangle_g$, which is $\overline F_g(\tilde q)$ by the definition of $\overline F_g$ (\cref{def:set-graphsum,def:set-modcohft}); grouping it by $(l,j)$ gives the right side. (e3) For $j\geq 1$, $\ML{g}$~(i) with $X_1 = \dots = X_j = H$ ($\Pi_0(H) = 1$) says that $\langle H^j\rangle_g-(\Fm_g)^{(j)}(v)$ lies in $\bbQ[v]$; both terms are homogeneous of exponent $2-2g-j<0$ (\cref{lem:set-modhom}; $\Fm_g = c_gv^{2-2g}$ by $\ML{g}$~(ii)), so the difference is $0$; for $j = 0$, $\ML{g}$~(ii) gives $\langle\ \rangle_g = c_gv^{2-2g}$. Hence
\[
\sum_j\frac{\delta^j}{j!}\langle H^j\rangle_g = c_gv^{2-2g}\sum_j\binom{2-2g}{j}\Bigl(\frac{\delta}{v}\Bigr)^j = c_g(v+\delta)^{2-2g}:
\]
$\delta/v$ lies in $u\bbQ\fps{u}$, as $v$ is $u$ times a unit and $\delta\in u^2\bbQ\fps{u}$, and the binomial identity $(1+x)^{2-2g} = \sum_j\binom{2-2g}{j}x^j$ of $\bbQ\fps{x}$ is preserved by the continuous substitution $x\mapsto\delta/v$. For $l\geq 1$, each term of the frame expansion is a correlator of homogeneous elements of $\sK$ containing at least one frame, on which $\Pi_0$ vanishes; by $\ML{g}$~(i) it lies in $\bbQ[v]$, so the term lies in $\bbQ\fps{u}$, and so does the sum of the summable family of these terms.

(f) The dilaton equation (\cref{lem:set-dilaton}~(i), which applies to the model by \cref{lem:set-dilaton}~(ii), in the form $\langle z1\,Y_1\cdots Y_r\rangle_h = (2h-2+r)\langle Y_1\cdots Y_r\rangle_h$ for $2h-2+r>0$), applied $m_1$ times, gives $\langle X(z1)^{m_1}y^{m_2}\rangle_h = (N)_{m_1}\langle Xy^{m_2}\rangle_h$, where $X := X_1\cdots X_n$, $N := 2h-2+n+m_2\geq 1$ and $(N)_{m_1} := N(N+1)\cdots(N+m_1-1)$; every intermediate step has $2h-2+r>0$. By multilinearity and symmetry the $m$-th term on the left is $\sum_{m_1+m_2 = m}\frac{c^{m_1}}{m_1!\,m_2!}\langle X(z1)^{m_1}y^{m_2}\rangle_h$, and for fixed $m_2$ the series in $m_1$ is $\frac{1}{m_2!}\langle Xy^{m_2}\rangle_h$ times the binomial series $\sum_{m_1}\binom{N+m_1-1}{m_1}c^{m_1}$. (f1) If $\varkappa(0) = 1$, then $\ord_uc\geq 1$ and the binomial series converges to $(1-c)^{-N} = \varkappa^N$. The family is summable: $\langle Xy^{m_2}\rangle_h = 0$ for $m_2>3h-3+n$ by (d), $\ord_uc^{m_1}\geq m_1$, and the finitely many correlators $\langle Xy^{m_2}\rangle_h$ lie in one $u^{N_0}\bbQ\fps{u}$ by (b1). Regrouping by $m_2$ gives $\sum_{m_2}\varkappa^{2h-2+n+m_2}\langle Xy^{m_2}\rangle_h/m_2!$, which is the right side. (f2) Let $c_0 := c(0)\neq 0$, $A_{m_2} := \langle Xy^{m_2}\rangle_h/m_2!$, let $o$ be the least $u$-order of the non-zero $A_{m_2}$ (finitely many, by (d)), $S$ the set of $m_2$ with $\ord_uA_{m_2} = o$, $a_{m_2} := \coef{u^o}A_{m_2}$ and $N_{m_2} := 2h-2+n+m_2$. Since $\ord_uc\geq 0$ and $\coef{u^0}c^k = c_0^k$, for $m\geq\max S$ the coefficient of $u^o$ in the $m$-th term $T_m$ is
\[
\coef{u^o}T_m = c_0^m\sum_{m_2\in S}c_0^{-m_2}a_{m_2}\binom{N_{m_2}+m-m_2-1}{N_{m_2}-1}.
\]
Each binomial coefficient is a polynomial in $m$ of degree $N_{m_2}-1$, and these degrees are distinct; so the sum is a non-zero polynomial in $m$, and $\coef{u^o}T_m\neq 0$ for all large $m$. Thus $T_m$ does not tend to $0$, and the series has no $u$-adic limit.
\end{proof}

In the applications of (f) (\cref{lem:tra-F}), $\varkappa = d$ and $y = T_{\geq 2}$, so that $cz1+y = q_T+z1$ and $\varkappa y = \tilde q_T+z1$ by \eqref{eq:tra-leaf}.
